\documentclass[12pt,a4paper]{article}
\usepackage[utf8]{inputenc}
\usepackage[russian]{babel}
\usepackage{geometry}
\usepackage{hyperref}
\usepackage{titlesec}
\geometry{a4paper, margin=1in}
\hypersetup{colorlinks=true,linkcolor=black,citecolor=black,urlcolor=blue}
\titleformat*{\section}{\large\bfseries}
\title{CryptoDigest --- Выпуск №1 (пилот)}
\date{9 августа 2025}
\begin{document}
\maketitle
\noindent\textit{Пилотный выпуск демонстрирует формат будущих дайджестов. Часть событий и чисел иллюстративны.}

\section*{Вступительное слово}
Первая неделя августа показала высокий темп новостей: регуляторные сдвиги в США/ЕС, рост роли L2-экосистем, конверсия points в airdrop и новая волна раздач. Мы объединяем нарративный анализ с метриками Coherence Score (CS) и Stimulation Index (SI), чтобы заранее оценивать устойчивость и стимулирующий потенциал инфоповодов и проектов.

\section*{Основные темы недели}
\begin{itemize}
  \item Регуляторика США/ЕС; локальные налоговые изменения в APAC.
  \item L2 и масштабирование: релизы в Base и zk/L2-rollups; рост активности в Telegram-экосистеме.
  \item Экономика points: конкуренция за ликвидность пользователей.
  \item Листинги: средние CEX ускоряют mid-cap, крупные держат фильтры качества.
  \item Безопасность: инциденты на мелких DEX/бриджах усилили запрос на аудиты.
\end{itemize}

\section*{Ожидания vs реальность}
\begin{itemize}
  \item Ожидание: крупный листинг L2-токена подтянет ликвидность; Реальность: без всплеска объёмов из-за рассевания внимания на airdrop-кампании.
  \item Ожидание: массовый аирдроп игры создаст устойчивый ап-тренд; Реальность: краткий хайп, затем ротация; DAU ниже ожиданий.
  \item Ожидание: партнёрство с финтех-провайдером даст прорыв ончейн; Реальность: SI вырос краткосрочно, ончейн вернулся к тренду.
\end{itemize}

\section*{Листинги и аирдропы (8 проектов)}
CS/SI: 0--100 (оценки на основе открытых сигналов; пилотные).

\subsection*{1) Atlas Layer (ATLAS)}
Событие: листинг Bybit, Gate.io (спот), 13--15 авг.\\
Токеномика: FDV \$450M, циркул. 18\%, разлоки ежемес.\\
Нарратив: универсальный data-availability слой.\\
Оценки: CS 68, SI 72; Горизонт: 14 дней.\\
Тезис: дев-ритм и партнёрства; риск --- разлоки.

\subsection*{2) NeonX (NEON)}
Событие: аирдроп тестнет-участникам; клейм 10--17 авг.\\
Оценки: CS 61, SI 79; Горизонт: 7 дней.\\
Тезис: сильный хайп, продуктовая глубина не доказана.

\subsection*{3) OrbitLend (ORBL)}
Событие: листинг MEXC, Bitget; LP-фарминг.\\
Оценки: CS 70, SI 65; Горизонт: 30 дней.\\
Тезис: взвешенная модель риска, следим за качеством TVL.

\subsection*{4) ZKStream (ZKS)}
Событие: листинг OKX (спот+перпеты), 14 авг.\\
Оценки: CS 66, SI 74; Горизонт: 14 дней.\\
Тезис: технарратив привлекателен; ключ --- удержание.

\subsection*{5) BaseSwap (BSWP)}
Событие: аирдроп на Base, клейм 9--16 авг.\\
Оценки: CS 64, SI 71; Горизонт: 7 дней.\\
Тезис: хорошая UX, высокая конкуренция.

\subsection*{6) Solara (SOLA)}
Событие: листинг Bybit (инновац. зона), 12 авг.\\
Оценки: CS 62, SI 69; Горизонт: 14 дней.\\
Тезис: сильный маркетинг; важны retention и DAU.

\subsection*{7) LayerPort (LPORT)}
Событие: кросс-экосистемный аирдроп (Ethereum/L2).\\
Оценки: CS 71, SI 67; Горизонт: 30 дней.\\
Тезис: утилити-кейс; безопасность бриджа критична.

\subsection*{8) EigenX (EIGX)}
Событие: листинг Gate.io, KuCoin; стейкинг валидаторов.\\
Оценки: CS 73, SI 63; Горизонт: 30 дней.\\
Тезис: фундаментально сильный кейс; риск --- сложность UX.

\section*{Методика CS/SI (пилот)}
CS: согласованность заявлений с фактами, ритм разработки, безопасность, соответствие трендам.\\
SI: охваты и вовлечённость, рыночная реакция, новизна нарратива, скорость медиа-диффузии.\\
Агрегация: взвешенная сумма компонент с калибровкой на ретроспективе.

\section*{Ресурсы и сервисы для MVP (1--2 мес.)}
Данные/API: CoinGecko Pro/CMC, DefiLlama, GitHub API, X API, Telegram/Discord боты, News API.\\
LLM/эмбеддинги: облачные или локальные (аренда GPU).\\
Хранилища: Postgres/Timescale, объектное, векторное (Qdrant/PGVector).\\
Вычисления/оркестрация: 2--3 vCPU, Prefect/Airflow, FastAPI.\\
Веб: Next.js/Static + CDN; мониторинг: Grafana/Prometheus.\\
Команда: data engineer, ML/NLP, full-stack, продакт (0.5 FTE).\\
OPEX пилота: ~\$1.0--2.5k/мес без команды.

\section*{Дисклеймер}
Материал носит информационный характер и не является инвестиционной рекомендацией. Часть данных --- иллюстративные.

\end{document}
